\documentclass[12pt]{article}

\usepackage[a4paper,margin=1in]{geometry}
\usepackage[hidelinks]{hyperref}
\usepackage{enumitem}
\usepackage{longtable}
\usepackage{titlesec}
\usepackage{listings}
\usepackage{xcolor}

\title{CCValidator Documentation}
\author{CCValidator}
\date{\today}

\definecolor{codebg}{RGB}{248,248,248}
\definecolor{codestring}{RGB}{163,21,21}
\definecolor{codecomment}{RGB}{0,128,0}
\definecolor{codekeyword}{RGB}{0,0,180}

\lstdefinelanguage{CSharp}{
  morekeywords={abstract,as,base,bool,break,byte,case,catch,char,checked,class,const,continue,decimal,default,delegate,do,double,else,enum,event,explicit,extern,false,finally,fixed,float,for,foreach,goto,if,implicit,in,int,interface,internal,is,lock,long,namespace,new,null,object,operator,out,override,params,private,protected,public,readonly,ref,return,sbyte,sealed,short,sizeof,stackalloc,static,string,struct,switch,this,throw,true,try,typeof,uint,ulong,unchecked,unsafe,ushort,using,virtual,void,volatile,while,record,var,Task},
  sensitive=true,
  morecomment=[l]{//},
  morecomment=[s]{/*}{*/},
  morestring=[b]",
}

\lstset{
  language=CSharp,
  basicstyle=\ttfamily\small,
  keywordstyle=\color{codekeyword},
  stringstyle=\color{codestring},
  commentstyle=\color{codecomment},
  backgroundcolor=\color{codebg},
  frame=single,
  breaklines=true,
  tabsize=2,
  showstringspaces=false
}

\begin{document}
\maketitle
\tableofcontents
\newpage

\section{Overview}

CCValidator is a FluentValidation-inspired validation library targeting .NET 10 / C\# 14.

\subsection{Packages}
\begin{itemize}[noitemsep]
  \item \texttt{CCValidator} --- core engine (\texttt{AbstractValidator<T>}, rule building, execution)
  \item \texttt{CCValidator.DependencyInjection} --- DI registration helpers
  \item \texttt{CCValidator.AspNetCore} --- ASP.NET Core MVC model validation integration
  \item \texttt{CCValidator.Serilog} --- optional Serilog logger adapter
\end{itemize}

\section{Getting started}

\subsection{Define a validator}
\begin{lstlisting}
public sealed record Person(string? Name, int Age);

public sealed class PersonValidator : AbstractValidator<Person>
{
  public PersonValidator()
  {
    RuleFor(x => x.Name)
      .NotEmpty()
      .MaximumLength(50);

    RuleFor(x => x.Age)
      .InclusiveBetween(0, 150);
  }
}
\end{lstlisting}

\subsection{Validate}
\begin{lstlisting}
var validator = new PersonValidator();
var result = validator.Validate(new Person(Name: "", Age: 200));

if (!result.IsValid)
{
  foreach (var failure in result.Errors)
    Console.WriteLine($"{failure.PropertyName}: {failure.ErrorMessage}");
}
\end{lstlisting}

\subsection{Async validation}
\begin{lstlisting}
RuleFor(x => x.Name)
  .MustAsync(async (name, ct) =>
  {
    await Task.Delay(10, ct);
    return name is not null;
  });
\end{lstlisting}

\section{API reference (short)}

\subsection{Core (CCValidator)}
\begin{itemize}[noitemsep]
  \item \texttt{AbstractValidator<T>} --- define rules via \texttt{RuleFor(...)} and execute via \texttt{Validate}/\texttt{ValidateAsync}.
  \item \texttt{IValidator<T>} --- validator interface (used by DI and ASP.NET integration).
  \item \texttt{ValidationResult} --- contains \texttt{IsValid} and \texttt{Errors}.
  \item \texttt{ValidationFailure} --- contains \texttt{PropertyName}, \texttt{ErrorMessage}, and optional metadata.
  \item \texttt{ValidationContext<T>} --- selects rule sets.
  \item \texttt{CCValidatorOptions} --- defaults and policies (messages, cascade, exception behavior, logging).
\end{itemize}

\subsection{Rule building}
\begin{itemize}[noitemsep]
  \item \texttt{IRuleBuilderInitial<T,TProperty>} / \texttt{IRuleBuilderOptions<T,TProperty>} --- fluent validator methods.
  \item Customization: \texttt{WithMessage(...)} and \texttt{WithErrorCode(...)}.
\end{itemize}

\subsection{Dependency injection (CCValidator.DependencyInjection)}
\begin{itemize}[noitemsep]
  \item \texttt{AddValidatorsFromAssembly(...)} and \texttt{AddValidatorsFromAssemblyContaining<T>(...)}
  \item \texttt{AddCCValidator(...)} overloads for registering/configuring \texttt{CCValidatorOptions}
\end{itemize}

\subsection{ASP.NET Core (CCValidator.AspNetCore)}
\begin{itemize}[noitemsep]
  \item \texttt{AddCCValidatorAutoValidation()} --- inserts a model validator provider.
\end{itemize}

\subsection{Logging (CCValidator + CCValidator.Serilog)}
\begin{itemize}[noitemsep]
  \item \texttt{IValidationLogger} --- receives internal validation errors.
  \item \texttt{SerilogValidationLogger} --- optional adapter that logs to Serilog.
\end{itemize}

\section{Dependency injection}

\subsection{Register validators}
\begin{lstlisting}
using CCValidator.DependencyInjection;
using Microsoft.Extensions.DependencyInjection;

var services = new ServiceCollection();
services.AddValidatorsFromAssemblyContaining<MyValidator>();
\end{lstlisting}

Validators are registered as \texttt{Scoped} by default. You can override the lifetime:

\begin{lstlisting}
services.AddValidatorsFromAssemblyContaining<MyValidator>(ServiceLifetime.Singleton);
\end{lstlisting}

\subsection{Configure defaults via CCValidatorOptions}
\begin{lstlisting}
services.AddCCValidator(o =>
{
  o.DefaultCascadeMode = CascadeMode.Stop;
  o.ExceptionBehavior = ValidationExceptionBehavior.ConvertToFailure;
  o.InternalErrorCode = "CCV_INTERNAL";
});
\end{lstlisting}

\section{ASP.NET Core MVC integration}

\subsection{Register auto-validation}
\begin{lstlisting}
using CCValidator.AspNetCore;
using CCValidator.DependencyInjection;

var builder = WebApplication.CreateBuilder(args);

builder.Services
  .AddControllers()
  .AddCCValidatorAutoValidation();

builder.Services.AddValidatorsFromAssemblyContaining<MyValidator>();
\end{lstlisting}

\subsection{Notes}
\begin{itemize}[noitemsep]
  \item MVC model validation is synchronous; async rules are not executed during model binding.
  \item If no \texttt{IValidator<TModel>} is registered for a model type, no errors are produced.
\end{itemize}

\section{Localization}

Default messages come from \texttt{IValidationMessageProvider}. You can supply a custom provider, including a ResourceManager-backed provider.

\subsection{ResourceManager provider}
\begin{lstlisting}
using System.Resources;

var rm = new ResourceManager("My.Messages", typeof(MyValidator).Assembly);
var provider = new ResourceManagerValidationMessageProvider(rm);

var options = new CCValidatorOptions { MessageProvider = provider };
\end{lstlisting}

\section{Logging}

Internal exceptions during validation can be logged via \texttt{IValidationLogger}.

\subsection{Serilog adapter}
\begin{lstlisting}
using CCValidator.Serilog;
using Serilog;

var log = new LoggerConfiguration()
  .WriteTo.Console()
  .CreateLogger();

var options = new CCValidatorOptions
{
  Logger = new SerilogValidationLogger(log),
};
\end{lstlisting}

\section{Migration from FluentValidation}

\subsection{What maps directly}
\begin{itemize}[noitemsep]
  \item \texttt{AbstractValidator<T>} and \texttt{RuleFor(...)}
  \item Chaining validators (e.g. \texttt{NotEmpty().MaximumLength(50)})
  \item \texttt{WithMessage(...)} and \texttt{WithErrorCode(...)}
  \item \texttt{When}/\texttt{Unless}, rule sets
  \item \texttt{Validate}/\texttt{ValidateAsync}
\end{itemize}

\subsection{Key differences}
\begin{itemize}[noitemsep]
  \item Configuration is explicit via \texttt{CCValidatorOptions} (and DI registration), not global static state.
  \item The API surface is intentionally smaller than FluentValidation (see \texttt{COMPATIBILITY.md}).
  \item MVC integration runs synchronously (async rules do not run during model binding).
\end{itemize}

\end{document}
